\subsection{Mô tả hệ thống đề xuất}

Dựa trên việc khảo sát hệ thống bán sách trực tuyến của Fahasa, nhóm đề xuất xây dựng một hệ thống quản lý nhà sách nhằm hỗ trợ quản lý toàn diện các hoạt động kinh doanh, lưu thông vật phẩm, chăm sóc khách hàng thành viên và điều phối nhân sự trong nhà sách. Hệ thống được xây dựng theo định hướng của một mô hình nhà sách hiện đại, không chỉ hỗ trợ bán các loại sách, báo và văn phòng phẩm mà còn hỗ trợ các nghiệp vụ mở rộng như mua trực tiếp tại cửa hàng, giao hàng tận nhà, mượn vật phẩm, thuê vật phẩm, quản lý bản sao của sách, gia hạn thẻ thành viên và áp dụng các đợt giảm giá. Các nhóm người dùng chính của hệ thống bao gồm khách hàng, nhân viên đứng quầy, tài xế chuyển hàng và quản trị viên hệ thống. Trong đó, khách hàng là đối tượng trực tiếp thực hiện các giao dịch mua, mượn hoặc thuê vật phẩm; nhân viên đứng quầy phụ trách xử lý các giao dịch tại quầy; tài xế chuyển hàng đảm nhận việc vận chuyển các đơn giao tận nhà; còn quản trị viên có trách nhiệm cập nhật dữ liệu hàng hóa, kiểm soát tồn kho và hỗ trợ vận hành toàn hệ thống.

Trong hệ thống được đề xuất, thực thể \textit{NhanVien} đại diện cho toàn bộ nhân sự làm việc trong nhà sách. Mỗi nhân viên được định danh duy nhất bởi thuộc tính \textit{IDNhanVien} và được mô tả thông qua các thuộc tính như \textit{CMND\_CCCD}, \textit{CongViec}, \textit{GioiTinh}, \textit{NgaySinh}, \textit{SDT} và thuộc tính tổ hợp \textit{HoVaTen}, trong đó \textit{HoVaTen} được cấu thành từ hai thuộc tính thành phần là \textit{Ho} và \textit{Ten}. Bên cạnh đó, nhân viên còn tham gia vào mối liên kết \textit{QuanLy}, dùng để mô tả quan hệ quản lý nội bộ trong nhà sách, trong đó có thuộc tính \textit{BoPhanQuanLy} nhằm xác định bộ phận hoặc phạm vi quản lý tương ứng. Ngoài ra, mỗi nhân viên có thể có một hoặc nhiều người phụ thuộc được lưu trữ trong thực thể \textit{NguoiPhuThuoc}. Mỗi thực thể \textit{NguoiPhuThuoc} được nhận diện bởi thuộc tính \textit{CMND\_CCCD} và được mô tả thêm bởi thuộc tính tổ hợp \textit{HoVaTen} gồm \textit{Ho} và \textit{Ten}, cùng với thuộc tính \textit{QuanHe} để xác định mối quan hệ giữa người phụ thuộc và nhân viên.

Thực thể \textit{NhanVien} được chuyên biệt hóa thành ba lớp con là \textit{TaiXeChuyenHang}, \textit{NhanVienDungQuay} và \textit{Admin}. Lớp con \textit{TaiXeChuyenHang} đại diện cho các nhân viên chịu trách nhiệm vận chuyển đơn hàng đến khách hàng và được mô tả bởi các thuộc tính riêng là \textit{IDBangLaiXe} . Lớp con \textit{NhanVienDungQuay} đại diện cho các nhân viên trực tiếp phục vụ khách tại quầy và có thuộc tính riêng là \textit{NamKinhNghiem}. Lớp con \textit{Admin} đại diện cho các nhân viên quản trị hệ thống, chịu trách nhiệm cập nhật dữ liệu hàng hóa, theo dõi số lượng nhập vào và kiểm soát việc thay đổi dữ liệu trong kho có thuộc tính cấp bậc. Các \textit{Admin} tham gia vào mối liên kết \textit{CapNhat} với thực thể \textit{VatPham}, trong đó mối liên kết này mang các thuộc tính \textit{SoLuongNhapVao} và \textit{NgayCapNhat} để ghi nhận mỗi lần cập nhật dữ liệu vật phẩm.

Thực thể \textit{KhachHang} đại diện cho các đối tượng sử dụng dịch vụ của nhà sách và là một kiểu hợp (union type/category) được hình thành từ hai kiểu thực thể là \textit{CaNhan} và \textit{TapTheCongTy}. Mỗi \textit{KhachHang} được định danh bởi thuộc tính \textit{IDKhachHang} và có các thuộc tính như \textit{SDT} và \textit{DiemTichLuy}. Thực thể \textit{CaNhan} đại diện cho khách hàng cá nhân, được nhận diện bởi thuộc tính \textit{CCCD\_CMND}, có thuộc tính tổ hợp \textit{HoVaTen} gồm \textit{Ho} và \textit{Ten}, đồng thời có thêm thuộc tính đa  \textit{DiaChi}. Thực thể \textit{TapTheCongTy} đại diện cho khách hàng là tổ chức hoặc doanh nghiệp, được định danh bởi thuộc tính \textit{MaSoThue} và được mô tả thêm bởi thuộc tính \textit{TenCongTyTapThe}. Việc mô hình hóa \textit{KhachHang} dưới dạng category cho phép hệ thống quản lý thống nhất cả khách hàng cá nhân lẫn khách hàng tổ chức nhưng vẫn giữ được các đặc trưng dữ liệu riêng của từng nhóm.

Để hỗ trợ chương trình khách hàng thân thiết, hệ thống sử dụng thực thể \textit{TheThanhVien}. Mỗi thẻ thành viên được nhận diện bởi thuộc tính \textit{MaThe} và được mô tả bởi các thuộc tính \textit{NgayLapThe} và \textit{NgayHetHan}. Trong mô hình, thẻ thành viên được gắn với khách hàng cá nhân thông qua mối liên kết \textit{TaoBoi}, phản ánh việc thẻ được khởi tạo dựa trên thông tin của một khách hàng cụ thể. Ngoài ra, thẻ thành viên còn có thể được gia hạn thông qua mối liên kết \textit{GiaHanThe} với thực thể \textit{Admin}, trong đó mối liên kết này có thuộc tính \textit{ThoiHanMoi} nhằm lưu thời điểm hết hạn mới sau khi gia hạn.

Thực thể \textit{VatPham} đại diện cho các mặt hàng được lưu hành trong nhà sách, bao gồm các sản phẩm phục vụ mua bán, thuê và mượn. Mỗi vật phẩm được định danh bởi thuộc tính \textit{IDVatPham} và được mô tả bởi các thuộc tính như \textit{TenVatPham}, \textit{GiaNiemYet} và \textit{uongKhaDung}. Trong hệ thống, \textit{VatPham} được chuyên biệt hóa thành ba lớp con là \textit{Sach}, \textit{Bao} và \textit{VanPhongPham}. Lớp con \textit{Sach} dùng để quản lý các mặt hàng sách và được mô tả bởi các thuộc tính riêng như \textit{NamXB}, \textit{TacGia} và \textit{NXB}. Lớp con \textit{Bao} dùng để quản lý báo hoặc tạp chí, với các thuộc tính như \textit{TapChi}, \textit{So} và \textit{ToaSoan}. Lớp con \textit{VanPhongPham} dùng để quản lý các mặt hàng văn phòng phẩm, có các thuộc tính như \textit{LoaiHang} và \textit{NSX}. Ngoài ra, một số đầu sách còn được quản lý theo từng bản sao vật lý thông qua thực thể \textit{BanCopy}. Mỗi \textit{BanCopy} được định danh bởi thuộc tính \textit{MaSoBanCopy}, đồng thời có thuộc tính \textit{TinhTrang} để xác định tình trạng bản copy hiện tại của thực thể \textit{Sach} thông qua mối liên kết \textit{Cua}, cho phép hệ thống quản lý các bản sao cụ thể của từng đầu sách.

Để hỗ trợ hoạt động khuyến mãi, hệ thống sử dụng thuộc tính đa trị, phức hợp \textit{DotGiamGia} ở thực thể \textit{VatPham}. Mỗi đợt giảm giá được nhận diện bởi thuộc tính \textit{STT} và được mô tả bằng các thuộc tính \textit{NgayBatDau}, \textit{NgayKetThuc} và \textit{PhanTramGiam}. Thực thể này được sử dụng để lưu trữ các chương trình ưu đãi áp dụng trong từng khoảng thời gian nhất định cho các mặt hàng hoặc nhóm mặt hàng, phục vụ cho các chiến dịch bán hàng và kích cầu mua sắm trong nhà sách.

Các giao dịch mua hàng của khách hàng được lưu trữ trong thực thể \textit{DonMuaHang}. Mỗi đơn mua hàng được định danh bởi thuộc tính \textit{ID} và được mô tả bằng các thuộc tính như \textit{LoaiHoaDon}, \textit{NgayMua} và thuộc tính dẫn xuất \textit{GiaTriTongDon}. Một khách hàng có thể tạo nhiều đơn mua hàng khác nhau thông qua mối liên kết \textit{Mua}. Bên cạnh đó, mỗi đơn mua hàng có thể chứa nhiều vật phẩm, và mỗi vật phẩm cũng có thể xuất hiện trong nhiều đơn mua hàng khác nhau. Mối quan hệ nhiều--nhiều này được biểu diễn thông qua mối liên kết \textit{Co} giữa \textit{DonMuaHang} và \textit{VatPham}, trong đó mối liên kết \textit{Co} mang các thuộc tính \textit{SoLuong} và \textit{GiaLucMua} để phản ánh số lượng từng vật phẩm trong đơn cũng như đơn giá thực tế tại thời điểm mua.

Thực thể \textit{DonMuaHang} tiếp tục được chuyên biệt hóa thành hai lớp con là \textit{MuaTrucTiep} và \textit{GiaoTanNha}. Lớp con \textit{MuaTrucTiep} đại diện cho các đơn hàng được thanh toán và nhận trực tiếp tại cửa hàng. Các đơn này được xử lý bởi nhân viên đứng quầy thông qua mối liên kết \textit{PhuTrach}, phản ánh trách nhiệm của từng nhân viên trong quá trình phục vụ và hoàn tất giao dịch tại quầy. Lớp con \textit{GiaoTanNha} đại diện cho các đơn hàng được giao đến địa chỉ của khách hàng và có các thuộc tính riêng như \textit{DiaChiNhanHang}, \textit{TrangThaiDon} và \textit{PhiVanChuyen}. Các đơn giao tận nhà được thực hiện thông qua mối liên kết \textit{ChuyenHang} với thực thể \textit{TaiXeChuyenHang}, cho phép hệ thống xác định nhân viên giao hàng phụ trách từng đơn cũng như theo dõi quá trình vận chuyển.

Ngoài chức năng mua hàng, hệ thống còn hỗ trợ hai loại hình dịch vụ lưu thông vật phẩm là mượn và thuê. Đối với hoạt động mượn, hệ thống sử dụng thực thể yếu \textit{LanMuon} để lưu thông tin cho từng lần mượn. Mỗi \textit{LanMuon} được mô tả bởi các thuộc tính như \textit{MaDonMuon}, \textit{NgayMuon}, \textit{HanTra}, \textit{NgayTra} và \textit{TinhTrangSauKhiTra}. Hoạt động mượn được hình thành thông qua mối liên kết định danh \textit{DungMuon}, trong đó khách hàng sử dụng thẻ thành viên để thực hiện việc mượn vật phẩm theo các quy định của hệ thống. Đối với hoạt động thuê, hệ thống sử dụng thực thể yếu \textit{LanThue} để lưu thông tin cho từng lần thuê. Mỗi \textit{LanThue} được mô tả bởi các thuộc tính như \textit{MaDonThue}, \textit{NgayThue}, \textit{ThoiHanTra}, \textit{NgayTra}, \textit{GiaThue} và \textit{SoLuong}. Hoạt động thuê được gắn với khách hàng và vật phẩm thông qua mối liên kết \textit{ThueVatPham}, cho phép hệ thống theo dõi các giao dịch thuê và quản lý việc hoàn trả vật phẩm sau khi sử dụng.

Nhờ cách tổ chức các kiểu thực thể, các thuộc tính và các mối liên kết như trên, hệ thống có thể phản ánh tương đối đầy đủ các nghiệp vụ cốt lõi của một nhà sách hiện đại lấy cảm hứng từ Fahasa, bao gồm quản lý nhân sự, quản lý khách hàng, quản lý vật phẩm, kiểm soát bán hàng, giao hàng, khuyến mãi, thẻ thành viên và các hoạt động mượn hoặc thuê vật phẩm. Mô hình dữ liệu này đồng thời tạo nền tảng phù hợp để tiếp tục xây dựng sơ đồ EERD và ánh xạ sang lược đồ cơ sở dữ liệu quan hệ trong các bước tiếp theo của bài tập lớn.

\subsection{Mô tả các ràng buộc ngữ nghĩa không biểu diễn đầy đủ bằng EERD}

\begin{enumerate}
    \item Khi khách hàng hoàn trả vật phẩm trong một thực thể \textit{LanMuon}, nếu thuộc tính \textit{TinhTrangSauKhiTra} được ghi nhận là hư hỏng hoặc mất mát, hệ thống phải tự động xác định mức bồi thường dựa trên thuộc tính \textit{GiaNiemYet} của vật phẩm tương ứng. Mức bồi thường được quy định nằm trong khoảng từ 100\% đến 150\% giá niêm yết tùy theo mức độ hư hỏng và loại vật phẩm.

    \item Hệ thống không cho phép một người dùng thuộc lớp con \textit{Admin} thực hiện phê duyệt, tạo mới hoặc gia hạn các giao dịch mượn, thuê hoặc thẻ thành viên cho các khách hàng là người thân của chính mình nếu những khách hàng đó xuất hiện trong mối liên kết với thực thể \textit{NguoiPhuThuoc}. Ràng buộc này nhằm đảm bảo tính khách quan và tránh xung đột lợi ích trong quá trình vận hành hệ thống.

    \item Trong vòng 30 ngày kể từ thuộc tính \textit{NgayLapThe} của thực thể \textit{TheThanhVien}, khách hàng chỉ được phép tồn tại tối đa hai giao dịch mượn còn hiệu lực cùng lúc. Sau khi vượt qua giai đoạn này, hệ thống mới cho phép khách hàng sử dụng đầy đủ hạn mức mượn theo chính sách thành viên.

    \item Đối với người dùng thuộc lớp con \textit{NhanVienDungQuay}, hệ thống chỉ cho phép hiển thị ba chữ số cuối của thuộc tính \textit{SDT} trong thực thể \textit{KhachHang}; các chữ số còn lại phải được che bằng ký tự thay thế. Trong khi đó, người dùng thuộc lớp con \textit{Admin} mới có quyền xem đầy đủ thông tin liên lạc của khách hàng.

    \item Hệ thống tự động khóa chức năng tạo mới giao dịch mượn hoặc thuê đối với những vật phẩm có số lượng khả dụng thấp hơn hoặc bằng 3 đơn vị. Quy định này nhằm ưu tiên các vật phẩm còn ít tồn kho cho hoạt động bán hàng thay vì các dịch vụ mượn hoặc thuê.

    \item Khách hàng chỉ được phép thực hiện giao dịch mượn hoặc thuê khi thẻ thành viên tương ứng còn hiệu lực, tức là thuộc tính \textit{NgayHetHan} của thực thể \textit{TheThanhVien} phải lớn hơn hoặc bằng ngày phát sinh giao dịch. Nếu thẻ đã hết hạn, hệ thống phải từ chối việc tạo mới \textit{LanMuon} hoặc \textit{LanThue}.

    \item Đối với các vật phẩm thuộc lớp con \textit{Sach} hoặc \textit{VanPhongPham} có thuộc tính \textit{GiaNiemYet} lớn hơn 2{,}000{,}000 VNĐ, khách hàng khi thực hiện giao dịch thuê phải thanh toán trước một khoản đặt cọc tối thiểu bằng 50\% giá trị niêm yết của vật phẩm trước khi giao dịch \textit{LanThue} được xác nhận.

    \item Hệ thống áp dụng hạn mức mượn khác nhau cho từng loại khách hàng. Cụ thể, khách hàng thuộc thực thể \textit{CaNhan} và khách hàng thuộc thực thể \textit{TapTheCongTy} có số lượng vật phẩm được phép mượn cùng lúc khác nhau theo chính sách của nhà sách. Ràng buộc này được sử dụng để kiểm soát lưu thông vật phẩm và giảm thiểu rủi ro thất thoát.

    \item Hệ thống tự động nâng cấp trạng thái ưu tiên của khách hàng khi tổng giá trị các đơn mua hàng của khách hàng trong một năm tài chính đạt từ 5{,}000{,}000 VNĐ trở lên. Khi đó, thuộc tính \textit{DiemTichLuy} hoặc trạng thái phân hạng khách hàng tương ứng sẽ được cập nhật để phản ánh quyền lợi thành viên nâng cao.

    \item Nếu một khách hàng có từ ba lần trả trễ hạn trở lên trong vòng sáu tháng, được xác định bằng việc so sánh thuộc tính \textit{NgayTra} với \textit{HanTra} trong các thực thể \textit{LanMuon}, hệ thống phải tự động tạm khóa quyền tạo mới giao dịch mượn hoặc thuê trong thời hạn 30 ngày kể từ thời điểm phát hiện vi phạm gần nhất.
\end{enumerate}